\chapter*{Introdução}

Hoje, é comum que se ouça, das mais diversas fontes, a alegação de que ``o celular nos escuta''.
Há aspectos corretos e errados nessa frase.
Ela é errada, no sentido de que esse tipo de vigilância não acontece literalmente escutando, com microfones e grampos, principalmente devido à escandalosidade desse fato e a dificuldade no tratamento de gravações de voz.
Ela acerta, contudo, na existência de uma coleta de dados por meios digitais na alma de muitos grandes negócios da Internet, incluindo aí as redes sociais -- Facebook, Instagram, TikTok, entre outros -- porém ocorrendo em termos muito mais discretos do que enuncia.

Dados aparentemente inofensivos e de produção rotineira, como históricos de visualização, \textit{likes}, trajetos de mobilidade, ingressos, compras, buscas, entre outros, são congregados e processados por algoritmos de inteligência artificial.
Com isso, cria-se um perfil comportamental do usuário, um modelo matemático que almeja prever suas ações.
\textcite{youyou_computerbased_2015} demonstraram que, tomando como base informações aparentemente inofensivas -- amizades do \textit{Facebook} -- é possível criar modelos computacionais que fazem previsões melhores, ou seja, sabem mais sobre o comportamento de uma pessoa, do que parentes e amigos próximos.
De modo similar, infere-se posicionamento político, localização, hábitos de consumo, interesses midiáticos e outra multitude de características pessoais extremamente valiosas partindo desses dados.
Esses modelos são fundamentais para as plataformas digitais, permitindo, por exemplo, que o anúncio mais eficaz seja remetido à pessoa; pode-se, assim, racionalizar e otimizar o comportamento humano a fim de maximizar seu engajamento, atenção, distração, desejo, tempo gasto e, enfim, o lucro da plataforma.

Essa monografia consiste no trabalho de tentar entender \textit{que tipo} de sociedade é criada pela coleta e processamento de dados, que hoje incorpora serviços utilizado por milhões, e que fazem parte importante, por exemplo, na formação de opinião pública.
Não se trata apenas de coleta de dados, mas também de concentração desses dados nas mãos de poucas empresas multinacionais, a maioria estadunidense, que, como se argumentará, passam a concentrar também poder político e econômico nesses processos.

No primeiro capítulo, serão discutidos os aspectos técnicos da vigilância.
Como ela ocorre?
Quem está por trás dela?
Que poder essas organizações detêm?
Nele, são descritas as bases técnicas por trás da sociedade de vigilância.

Em seguida, no segundo capítulo, se discutem as inteligências artificiais, partes cruciais do processamento de dados pessoais das quais muito se fala, frequentemente sem o aprofundamento técnico necessário.
Através da pesquisa dos fundamentos matemáticos por trás do aprendizado de máquina, esses algoritmos são analisados em termos de transparência e interpretabilidade.
As escolhas técnicas de \textit{como} são feitos esses algoritmos tem impactos reais nos seus usuários, principalmente, pois esses se tratam de aproximações, e podem (e vão) errar.

O terceiro capítulo analisa o Vale do Silício, e como ele se consolidou como uma potência ideológica e política através da Internet.
O argumento principal é que essas indústrias se aproveitaram da onda neoliberal dos anos 80 e 90 e, por meio de um discurso ideológico específico, afastou por tempo suficiente a intervenção e regulação estatal do meio digital.
Com isso, pode se proliferar livremente e tomar controle sobre esse espaço.
Por fim, nesse mesmo capítulo é apresentado um modelo de produção e organização alternativo ao capitalismo do Vale do Silício -- o \textit{software} livre.
